\section{Related Work}
\label{sec:related}

Our work sits at the intersection of three research streams: LLM code evaluation, property-based testing for code verification, and formal contract checking.

\subsection{LLM Code Evaluation: From Sparse to Dense Testing}

The HumanEval benchmark \citep{Chen2021HumanEval} established functional correctness via unit tests as the primary evaluation axis for LLM-generated code. EvalPlus \citep{Liu2023EvalPlus} addressed test-gaming by expanding HumanEval by 80$\times$ and MBPP by 35$\times$, reducing pass@1 by up to 23.1\%. This dense differential testing approach represents the current gold standard for test-based LLM code evaluation.

However, differential testing is limited to checking \emph{output equality on sampled inputs}. It cannot check whether a program satisfies properties that must hold for all valid inputs: relational invariants, quantified conditions, or semantic contracts that specify behavior beyond input-output pairs. Our work builds on EvalPlus as the strongest available test-based baseline and quantifies how much weaker it is than a formal contract oracle.

\subsection{Property-Based Testing for LLM Code}

\citet{Bose2025Prompts} applied PBT (Hypothesis-style) to StarCoder and CodeLlama on MBPP and HumanEval, finding that 18--32\% of programs fail outright. However, this work covers only two models, uses no formal contract annotations, and lacks the oracle isolation design needed to separate input-generation advantage from oracle semantic strength.

\citet{Newcomb2025Preconditions} examined how pre/postcondition constraints provided as prompts affect LLM code generation quality, showing that structural constraints improve correctness. This complements our work: where they study contract-guided generation, we study contract-based evaluation of existing outputs.

\textbf{The key gap:} No prior PBT work answers whether additional violations reflect \emph{oracle semantic strength} or \emph{input exploration advantage}. Our oracle isolation design (Section~\ref{sec:method}) resolves this directly.

\subsection{Formal Contract Checking for LLM Code}

ContractEval \citep{Lim2025ContractEval} is the only publicly available benchmark pairing HumanEval+/MBPP+ tasks with formal Python pre/post-condition contracts. Their evaluation found 0\% contract satisfaction for 5 open-source models under standard prompting. Critically, ContractEval uses Z3 SMT checking (25.82\% tractable). Our execution-based approach achieves 100\% tractability across all 364 tasks and extends evaluation to closed-source models.

All claims in this related work survey are drawn from peer-reviewed or publicly archived sources; we do not rely on unverified external claims.

\subsection{Summary}

\begin{table}[h]
\centering
\small
\caption{Comparison with related work.}
\label{tab:related}
\begin{tabular}{lllll}
\toprule
Work & Oracle Type & Models & Oracle Isolation & Contract Annot. \\
\midrule
EvalPlus \citep{Liu2023EvalPlus} & Differential (dense) & Multiple & --- & None \\
\citet{Bose2025Prompts} & PBT (no formal) & 2 & No & Manual \\
ContractEval \citep{Lim2025ContractEval} & SMT (25.82\% tractable) & 5 (open) & No & Formal \\
\citet{Newcomb2025Preconditions} & Test-based & Multiple & No & Prompt-provided \\
\textbf{This work} & \textbf{Exec.-based (100\%)} & \textbf{5 (open+closed)} & \textbf{Yes (CVT)} & \textbf{Formal} \\
\bottomrule
\end{tabular}
\end{table}
